\documentclass[a4paper,8pt]{article}

\usepackage[top=0.25in, bottom=0.25in, left=0.4in, right=0.4in]{geometry}
\usepackage{fontspec}
\usepackage{enumitem}
\usepackage{titlesec}
\usepackage[hidelinks]{hyperref}

\setmainfont{Helvetica}
\setlength{\parindent}{0pt}
\setlength{\parskip}{1pt}

\titleformat{\section}{\normalsize\bfseries}{}{0em}{}[\titlerule]
\titlespacing*{\section}{0pt}{7pt}{3pt}

\setlist[itemize]{nosep, topsep=1pt, leftmargin=1.5em, label=\textbullet, itemsep=0pt, parsep=0pt}

\pagestyle{empty}

\begin{document}

% ===== KİŞİSEL BİLGİLER =====
\begin{center}
{\Large\bfseries Emre Pirinç}\\[2pt]
{\small Aspiring AI Engineer | Machine Learning | LLM | Generative AI | Python | Data Science | MIS Graduate}\\[2pt]
{\small İstanbul, Türkiye | info@1takimstartuplar.com | +90 534 590 58 89}\\[1pt]
{\small LinkedIn: \href{https://linkedin.com/in/emre-pirinc}{linkedin.com/in/emre-pirinc} \textbar{} GitHub: \href{https://github.com/EmrePirinc}{github.com/EmrePirinc}}
\end{center}

\vspace{-2pt}

% ===== HAKKIMDA =====
\section{HAKKIMDA}
{\small Sakarya Üniversitesi Yönetim Bilişim Sistemleri mezunuyum. Yapay zeka, makine öğrenmesi ve büyük dil modelleri (LLM) alanlarında aktif olarak projeler geliştiriyorum. Gemini Live API, Stable Diffusion ve çoklu AI model entegrasyonları (Gemini, GPT, Claude, DeepSeek) üzerine pratik deneyim sahibiyim.

Kariyerime Yapay Zeka Mühendisliği alanında devam etmek istiyorum. Arkadaşlarımla birlikte AI tabanlı 3 proje üzerinde çalışıyor ve startup yarışmalarına hazırlanıyoruz. Uzun vadeli hedefim; AI destekli kurumsal çözümler ve SAP Joule AI sistemleri üzerine uzmanlaşmak.}

% ===== YETENEKLER =====
\section{YETENEKLER}
{\small Python, Artificial Intelligence, Machine Learning, LLM, Generative AI, Prompt Engineering, Data Science, Data Analysis, Data Engineering, SQL, Gemini API, Stable Diffusion, Chrome Extension Development, WebSocket, Web Audio API, Natural Language Processing, SAP Business One, SAP ERP, ERP Implementation, Crystal Reports, Blueprint / Functional Design, Business Analysis, Requirements Analysis, Project Management, UAT, System Integration, Test Automation, Veritabanı Yönetimi, ISO 9001, Scrum, Figma, Kanban}

% ===== İŞ DENEYİMİ =====
\section{İŞ DENEYİMİ}

\textbf{1 Takım Start Uplar} \hfill {\small Şubat 2026 -- Devam Ediyor}\\
{\small \textit{Kendi İşim} -- İstanbul, Türkiye}
\begin{itemize}
\item {\small 5 kişilik ekiple 3 aktif proje geliştiriyoruz: DeepNote (AI destekli not uygulaması -- ses kaydı, transkript, özet, akış diyagramı), DeepLive (web sayfalarındaki yabancı dili gerçek zamanlı çeviren Chrome eklentisi -- Gemini Live API), DeepSlide (AI ses analizi ile anahtar kelimelere göre otomatik slayt geçişi yapan sunum uygulaması).}
\end{itemize}

\vspace{3pt}

\textbf{AIFTeam | SAP Business One Certified Partner} \hfill {\small Şubat 2025 -- Şubat 2026 · 1 yıl 1 ay}\\
{\small \textit{İstanbul, Türkiye}}

\vspace{3pt}
{\small \textbf{HERO Ürün Danışmanı} · Tam Zamanlı \hfill Ocak 2026 -- Şubat 2026}\\
{\small Uzaktan}
\begin{itemize}
\item {\small AIFTeam bünyesinde yürütülen Hero projesinde; HMR ve PM modülleri ile mobil uygulamaları içeren entegre portalın süreç yönetimi ve operasyonel verimlilik çalışmalarında aktif rol almaktayım.}
\end{itemize}

\vspace{3pt}
{\small \textbf{SAP B1 Consultant} · Tam Zamanlı \hfill Temmuz 2025 -- Şubat 2026}\\
{\small Ofisten}
\begin{itemize}
\item {\small SAP Business One danışmanlık süreçlerinde analiz, tasarım, geliştirme ve canlıya geçiş aşamalarının tüm süreç yönetiminde görev almaktayım.}
\item {\small Proje Yönetimi \& Analiz: Müşteri ihtiyaç analizlerinin yapılması, Kavramsal Tasarım Dokümanlarının (Blueprint) hazırlanması ve proje süreçlerinin takibi.}
\item {\small Raporlama \& Tasarım: 3 Crystal Reports raporunu sıfırdan tasarlayıp 10+ raporda düzeltme sağladım. 4 SAP raporunu SQL ile sıfırdan geliştirdim. Form dizaynları ve sorgu optimizasyonu yaptım.}
\item {\small Kurulum \& Konfigürasyon: Şirket, Server, Client ve SAP B1 kurulumlarının gerçekleştirilmesi ve sistem parametrelerinin uyarlanması.}
\item {\small Entegrasyon \& Geliştirme: Web projeleri ve 3. parti yazılımlar ile SAP B1 entegrasyonlarının yönetimi; add-on süreçlerine destek verilmesi.}
\item {\small Test \& Destek: Kullanıcı Kabul Testleri (UAT) ve canlı kullanım sonrası destek süreçlerinin yürütülmesi. Web ekranı geliştirmesi için 5 kişilik ekibe eğitim verildi.}
\end{itemize}

\vspace{3pt}
{\small \textbf{SAP B1 Consultant} · Stajyer \hfill Şubat 2025 -- Haziran 2025}\\
{\small Ofisten}
\begin{itemize}
\item {\small 6 farklı müşteride yerinde destek ve gözlem yaparak müşteri iletişimi, teknik spesifikasyon dokümanı hazırlama, rapor yazımı ve Crystal Reports kullanımı üzerine mesleki gelişim sağladım.}
\end{itemize}

\vspace{3pt}

\textbf{Intern} \hfill {\small Temmuz 2024 -- Eylül 2024}\\
{\small \textit{Seratu Medya Reklam Yazılım Danışmanlık · Stajyer} -- Konya, Türkiye}
\begin{itemize}
\item {\small Ortaokul seviyesine yönelik Bilişim ve Teknoloji Tasarım kitaplarının scratch robotik kodlama ve algoritma içeriklerini hazırladım. Görsel tasarım süreçlerinde Stable Diffusion kullanarak yapay zeka destekli içerikler ürettim ve QR kod entegrasyonu ile interaktif ders materyalleri tasarladım.}
\end{itemize}

\vspace{3pt}

\textbf{Paketleme} \hfill {\small Ağustos 2023 -- Ekim 2023}\\
{\small \textit{TTT-auto · Dönemsel} -- Konya, Türkiye}
\begin{itemize}
\item {\small Global ölçekte üretim yapan firmada depo yönetim sistemleri (WMS) ve otomasyon süreçlerini yerinde deneyimleyerek lojistik operasyonlarını anlamlandırdım.}
\end{itemize}

\vspace{3pt}

\textbf{TİİBÖT Sosyal Medya Editörü} \hfill {\small Mart 2022 -- Haziran 2022}\\
{\small \textit{TİİBÖT} -- Türkiye}
\begin{itemize}
\item {\small Öğrenci topluluğunun dijital iletişim süreçlerinde gönüllü olarak görev aldım. Etkinlikler için sosyal medya içerikleri hazırladım ve organizasyonlara uzaktan katılarak yönetim ekibine operasyonel destek verdim.}
\end{itemize}

% ===== PROJELER =====
\section{PROJELER}

\textbf{Anadolu Bakır SAP Business One Canlı Geçiş Projesi} \hfill {\small Temmuz 2025 -- Şubat 2026}\\
{\small AIFTeam | SAP Business One Certified Partner ile ilişkili.}

\vspace{3pt}

\textbf{QDMS Kalite Doküman Yönetimi Sistemi Tasarlanması} \hfill {\small Kasım 2024 -- Aralık 2024}\\
{\small 7 haftalık sprintler boyunca ekip olarak ISO 9001 standartlarına uygun bir kalite yönetim sistemi geliştirdik. Müşteri gereksinimlerini analiz ederek çalışan bir yazılım sunduk. Sakarya Üniversitesi ile ilişkili.}

\vspace{3pt}

\textbf{SMA'YA IŞIK TUT} \hfill {\small Ocak 2024 -- Haziran 2024}\\
{\small SMA ve bağışçıları bir araya getiren ve farkındalık amaçlayan bir mobil uygulamanın teorik çalışması. Gereksinim belirleme, UI/UX tasarımı, proje planlama ve risk yönetimi. Sakarya Üniversitesi ile ilişkili.}

\vspace{3pt}

\textbf{Uber Gereksinim Analizi Yapılması} \hfill {\small Ocak 2024 -- Haziran 2024}\\
{\small Uber iş modelini detaylı analiz ederek gereksinimlerini tanımladık. Use case diyagramları, wireframe, storyboard ve story mapping tekniklerinin kullanımı. Sakarya Üniversitesi ile ilişkili.}

\vspace{3pt}

\textbf{Sakarya Üniversitesi Öğrencilerinin Yapay Zeka Bağımlığı Ve Korkuları} \hfill {\small Kasım 2023 -- Devam Ediyor}\\
{\small TÜBİTAK 2209-A destekli proje. Sonuç raporu yükleme aşamasında. Sakarya Üniversitesi ile ilişkili.}

\vspace{3pt}

\textbf{E-Park: Akıllı Otopark Çözümleri YBS Projesi} \hfill {\small Ocak 2021 -- Haziran 2021}\\
{\small Yoğun trafik ve park sorunlarına çözüm getiren, mobil tabanlı, entegre akıllı otopark sistemi projesi. Plaka Tanıma Sistemi (PTS), Masterpass ödeme altyapısı ve Blockchain destekli veri güvenliği. Sakarya Üniversitesi ile ilişkili.}

% ===== EĞİTİM =====
\section{EĞİTİM}

\textbf{İstanbul Üniversitesi} \hfill {\small Ekim 2025 -- Temmuz 2027}\\
{\small \textit{Ön Lisans, Bilgisayar Programcılığı} -- Teknik programlama yetkinliğini güçlendirmek amacıyla}

\vspace{3pt}

\textbf{Sakarya Üniversitesi} \hfill {\small Ekim 2021 -- Haziran 2025}\\
{\small \textit{Lisans, Yönetim Bilişim Sistemleri} -- Not: 3.27}

\vspace{3pt}

\textbf{İzmir Bakırçay Üniversitesi} \hfill {\small Eylül 2019 -- Eylül 2021}\\
{\small \textit{Lisans, Yönetim Bilişim Sistemleri} -- İngilizce hazırlık eğitimi tamamlandı}

% ===== LİSANSLAR VE SERTİFİKALAR =====
\section{LİSANSLAR VE SERTİFİKALAR}
{\small
\textbf{Google Cloud Introduction to Data Engineering} -- Istanbul Data Science Academy (Mart 2026)\\
\textbf{CSCON'26 (IEEE Türkiye Computer Society Congress)} -- IEEE Computer Society Türkiye SAC (Mart 2026)\\
\textbf{Insider One AI Weekend Certification} -- Insider One (Şubat 2026)\\
\textbf{SAP Certified -- SAP Business One} -- SAP (Ocak 2026) -- \href{https://www.credly.com/badges/149da554-e1b6-477d-a063-f3c7997625c2/linked_in_profile}{Doğrula}\\
\textbf{Artificial Intelligence -- 2602 Early Release Series} -- SAP (Ocak 2026)\\
\textbf{Introduction to Large Language Models} -- Google Cloud Training Online (Ocak 2026)\\
\textbf{Veri Bilimi ve Yapay Zeka Zirvesi Datacamp'2025} -- Boğaziçi Üniversitesi (Aralık 2025)\\
\textbf{Python ile Programlamaya Giriş DATACAMP'25} -- Boğaziçi Üniversitesi (Aralık 2025)\\
\textbf{Querying MS SQL} -- Miuul (Şubat 2025)\\
\textbf{Uygulamalarla SQL Öğreniyorum} -- BTK Akademi (Ocak 2025)\\
\textbf{SAP Business One Katılım Sertifikası} -- UpexTech (Ocak 2025)\\
\textbf{Data Analysis in Excel} -- DataCamp (Kasım 2024)\\
\textbf{Time Series Analysis in Python} -- DataCamp (Kasım 2024)\\
\textbf{Veri Bilimi Zirvesi Datacamp24} -- Boğaziçi Üniversitesi (Kasım 2024)\\
\textbf{Business Analysis and PO with AI Assistance} -- BA-WORKS İş Analizi Hizmetleri (Ağustos 2024)\\
\textbf{Veri Ön İşleme} -- T.C. Cumhurbaşkanlığı Dijital Dönüşüm Ofisi (Mart 2024)\\
\textbf{Boğaziçi DataCamp'22 Veri Bilimi Zirvesi} -- Boğaziçi Üniversitesi (Kasım 2022)\\
\textbf{Boğaziçi DataCamp'21 Veri Bilimi Zirvesi} -- Boğaziçi Üniversitesi (Aralık 2021)
}

% ===== ONUR VE ÖDÜLLER =====
\section{ONUR VE ÖDÜLLER}
{\small
\textbf{Profesyonel E-Spor Kariyeri} -- (2020 -- 2021) -- 7 ulusal turnuvada 5 şampiyonluk, 1 Türkiye 2.'liği ve 1 profesyonel lig deneyimi. 2 turnuvada MVP ödülü. Tüm turnuvalarda IGL (In-Game Leader) olarak takım yönetimi, strateji belirleme ve kriz anlarında karar alma sorumluluğunu üstlendim.
}

% ===== DİLLER =====
\section{DİLLER}
{\small Türkçe (Ana Dil), English (A2)}

% ===== REFERANSLAR =====
\section{REFERANSLAR}
{\small
\textbf{Prof. Dr. Adem Akbıyık} -- Sakarya Üniversitesi, Yönetim Bilişim Sistemleri\\
adema@sakarya.edu.tr | +90 532 680 51 24\\[3pt]
\textbf{Prof. Dr. Erman Coskun} -- Bahçeşehir Üniversitesi, Yönetim Bilişim Sistemleri\\
coskune@gmail.com | +90 544 417 94 00
}

\end{document}
